\subsection{Database ব্যবস্থার প্রয়োগক্ষেত্র}
বিভিন্ন প্রতিষ্ঠানে কাজ আলাদা হলেও তথ্য সংরক্ষণ, সংযোগ, অনুসন্ধান ও প্রতিবেদন তৈরির প্রয়োজনটি সাধারণ। নিচের উদাহরণগুলোতে কোন ডেটা রাখা হয় এবং তা দিয়ে কী কাজ হয়, দুটোই লক্ষ করো।
\begin{multicols}{2}
\chSixLead{১}{সরকারি প্রশাসন} নাগরিক বা এলাকার পরিচিতির সঙ্গে সেবার তথ্য যুক্ত রাখা যায়। যেমন একটি পৌরসভা হোল্ডিং নম্বর দিয়ে কর পরিশোধের অবস্থা খুঁজে পায়।
\chSixLead{২}{ব্যাংকিং} গ্রাহক, হিসাব ও লেনদেনের পৃথক রেকর্ড মিলিয়ে স্থিতি ও হিসাববিবরণী তৈরি করা যায়। একই গ্রাহকের একাধিক হিসাব থাকলেও পরিচিতির মাধ্যমে তাকে চেনা যায়।
\chSixLead{৩}{ব্যবসা ও বিক্রয়} পণ্যের কোড, মজুদ, বিক্রয় ও ক্রেতার তথ্য থেকে বিল এবং অবশিষ্ট পণ্যের তালিকা তৈরি হয়। একই পণ্যের অনেক বিক্রয়রেকর্ড থাকতে পারে।
\chSixLead{৪}{পরিবহন ও বিমানসেবা} যাত্রা, সময়সূচি, যাত্রী ও আসন বরাদ্দ মিলিয়ে টিকিট দেওয়া হয়। একই যাত্রার একই আসন যেন দুজনকে বরাদ্দ না হয়, সে জন্য সামঞ্জস্যের বিধি দরকার।
\chSixLead{৫}{শিক্ষা} শিক্ষার্থীর ভর্তি, উপস্থিতি, ফি ও পরীক্ষার ফল আলাদা টেবিলে রেখেও পরিচিতির মাধ্যমে যুক্ত করা যায়। একটি কুয়েরি দিয়ে নির্দিষ্ট শ্রেণির ফল বের করা সম্ভব।
\chSixLead{৬}{টেলিযোগাযোগ} গ্রাহক পরিচিতি, ব্যবহার ও বিলের তথ্য মিলিয়ে মাসিক পাওনা নির্ণয় করা যায়। ভুল গ্রাহকের ব্যবহার অন্য হিসাবে যুক্ত না হওয়া জরুরি।
\chSixLead{৭}{আর্থিক ব্যবস্থাপনা} প্রতিষ্ঠানের আয়, ব্যয় ও সম্পদের তথ্য থেকে নির্দিষ্ট সময়ের সারাংশ তৈরি হয়। প্রতিটি হিসাবের উৎসরেকর্ড থাকায় ভুল শনাক্ত করা সহজ হয়।
\chSixLead{৮}{উৎপাদন} কাঁচামাল, উৎপাদিত পণ্য ও ক্রেতার চাহিদার তথ্য মিলিয়ে উৎপাদন পরিকল্পনা করা যায়। মজুদ কমে গেলে প্রয়োজনীয় উপকরণ শনাক্ত করা সম্ভব।
\chSixLead{৯}{মানবসম্পদ} কর্মীর পরিচিতি, উপস্থিতি ও বেতনের তথ্য থেকে মাসিক বেতনবিবরণী তৈরি হয়। একই নামের কর্মীদের পৃথক EmployeeID দিয়ে চেনা যায়।
\chSixLead{১০}{ওয়েবসেবা} ব্যবহারকারীর নিবন্ধন, জমা দেওয়া আবেদন বা মন্তব্য সংগঠিতভাবে রাখা যায়। অনুমতি অনুযায়ী নিজের আবেদন দেখার সুযোগ দেওয়া যায়।
\chSixLead{১১}{স্বাস্থ্যসেবা} রোগীর পরিচিতির সঙ্গে সাক্ষাৎ, পরীক্ষা ও চিকিৎসার রেকর্ড যুক্ত থাকে। একই রোগীর একাধিক সাক্ষাৎ আলাদা রেকর্ডে রেখে তাঁর ইতিহাস খোঁজা যায়।

\end{multicols}

